\documentclass{article}
\usepackage{amsmath , algorithmic ,natbib}
\title{Introduction to LaTeX}
\author{Bodike Chaithali \\ B24CS013}
\date{\today}
\begin{document}
	\maketitle
	\section*{Q1}
	\begin{abstract}
		An abstract reflecting on the role of LATEX in research writing \dots
	\end{abstract}
	\section{Introduction}
	This is an introductory section.
	\section{Methods}
	This is an methods section.
	\subsection{Data Collection}
	This is an Data Collection SubSection in Methods section.
	\subsection{Analysis}
	This is an Analysis SubSection in Methods section.
	\section{Conclusion}
	This is an conclusion section.
	\section*{Q2}
	This line refers to the "Key Formula" by the equation number as \dots \ref{eq:euler}\\
	\\
	Euler's equation is :
	\begin{equation}
		e^{i\pi} + 1 =0
		\label{eq:euler}
	\end{equation} 
	\section*{Q3}
	\begin{table}[h]
		\centering
	\begin{tabular}{|c|c|c|c|}
		\hline
		Student & Math & Physics & CS \\
		\hline
		Alice & 99 & 98 & 97 \\
		\hline
		Bob & 95 & 97 & 99 \\
		\hline
		Charlie & 99 & 93 & 100 \\
		\hline
	\end{tabular}
\end{table}
	\section*{Q4}
		\begin{algorithmic}
			\STATE Input: integer  $n$
			\IF{$n > 0$}
				\STATE print "Positive"
			\ELSIF{$n < 0$}
				\STATE print "Negative"
			\ELSE
				\STATE print "Zero"
			\ENDIF
		\end{algorithmic}
	\section*{Q5}
	An 1st one \citet{knuth1984} and the 2nd one \citep{lamport1994} . 
	\bibliographystyle{plainnat}
	\bibliography{Lab2}
\end{document}